\section{Proof of Lemma~\ref{lem:h}}

\begin{proof}
Suppose that the roots $\eta=\mathrm{i} s_j$, $j=1,2,3$ are distinct, and that the Boutroux conditions~\eqref{eq:Boutroux} hold (this will be justified later via a continuation argument). With an integration constant selected so that $\operatorname{Re}(h(\eta,y,c))=0$ at any one of the roots, the zero level set $K$ of $\operatorname{Re}(h(\eta,y,c))$ is well defined and it contains the closure $K'$ of the union of \emph{critical trajectories} emanating from the points $\eta=\mathrm{i} s_j$, $j=1,2,3$, which are the curves along which $f(\eta,y,c)\,\mathrm{d}\eta^2<0$, where $f$ is defined by~\eqref{eq:h-eta}. We claim that $K'$ has the following properties.
\begin{itemize}\itemsep=0pt
\item $K'$ is a connected set consisting of six simple arcs pairwise disjoint except for their endpoints:
\begin{itemize}\itemsep=0pt
\item One arc joining the origin $\eta=0$ to one of the three points that we label as $\eta=\mathrm{i} s_1$. We take this arc to be the branch cut $\Sigma_{0,1}$.
\item One arc joining the other two points, $\eta=\mathrm{i} s_2$ and $\eta=\mathrm{i} s_3$. We take this arc to be the branch cut $\Sigma_{2,3}$.
\item Two arcs joining $\eta=\mathrm{i} s_1$ either to the same point that we label as $\eta=\mathrm{i} s_2$ (case (i)) or one each to $\eta=\mathrm{i} s_2$ and $\eta=\mathrm{i} s_3$ (case (ii)).
\item Two unbounded arcs tending to $\eta=\infty$ parallel to the real line, one in the right half-plane and one in the left half-plane.
\end{itemize}
\item In case (i), the region bounded by the two arcs joining $\eta=\mathrm{i} s_1$ with $\eta=\mathrm{i} s_2$ contains the origin $\eta=0$ and both unbounded arcs emanate from $\eta=\mathrm{i} s_3$. In case (ii), the region bounded by the arc joining $\eta=\mathrm{i} s_1$ with $\eta=\mathrm{i} s_2$, the arc joining $\eta=\mathrm{i} s_1$ with $\eta=\mathrm{i} s_3$, and the arc $\Sigma_{2,3}$ contains the origin $\eta=0$ and one unbounded arc emanates from each of $\eta=\mathrm{i} s_2$ and $\eta=\mathrm{i} s_3$.
\end{itemize}
Locally, near each of the simple roots $\eta=\mathrm{i} s_j$, $j=1,2,3$, of $\eta\mapsto f(\eta,y,c)$, $K'$ consists of a union of three trajectories emanating from $\eta=\mathrm{i} s_j$ in directions separated by equal angles of $2\pi/3$. Given an index $j=1,2,3$, each of the three trajectories emanating from $\eta=\mathrm{i} s_j$ terminates in the other direction at $\eta=\mathrm{i} s_1$, $\eta=\mathrm{i} s_2$, $\eta=\mathrm{i} s_3$, $\eta=0$, or is unbounded in which case it tends to $\eta=\infty$ asymptotically horizontally. This is because otherwise the trajectory would be \emph{divergent} and hence \emph{recurrent}~\cite[Theorem 11.1]{Strebel84}. But the closure of a~recurrent trajectory contains a~nonempty domain in $\mathbb{C}$ and since $\operatorname{Re}(h(\eta,y,c))=0$ on the trajectory, this harmonic function would vanish identically on $\mathcal{R}$ (the Riemann surface of $h_\eta(\eta,y,c)$, i.e., the spectral curve), which is a contradiction because $h_\eta(\eta,y,c)$ is not identically zero. Taking into account the Boutroux conditions~\eqref{eq:Boutroux} which imply that $K'\subset K$, similar local analysis shows that there can be at most one critical trajectory terminating at $\eta=0$ and at most one unbounded critical trajectory tending horizontally to $\eta=\infty$ in each of the left and right half-planes.

To work out the global trajectory structure in order to prove the claim, it is easiest to first assume that $y>0$ and $c\in\mathbb{R}$, in which case it is easy to see that $h_\eta(-\eta^*,y,c)=h_\eta(\eta,y,c)^*$ provided that $\Sigma_{0,1}$ and $\Sigma_{2,3}$ are taken to be symmetric in the imaginary $\eta$-axis, which we will also assume. Moreover we either have (making a choice of labeling of the roots of $P$) $s_1<0<s_2<s_3$ or $s_1<0$ and $s_3=s_2^*$. In either configuration the condition $I_{2,3}=0$ holds automatically, and $c\in\mathbb{R}$ is presumed to be determined from the remaining real condition $I_{1,2}=0$. We examine the two configurations in turn.

If $s_1<0<s_2<s_3$, since $f(\eta,y,c)=h_\eta(\eta,y,c)^2>0$ holds for $\eta\in\mathrm{i}\mathbb{R}$ between $\eta=\mathrm{i} s_1$ and $\eta=0$ as well as between $\eta=\mathrm{i} s_2$ and $\eta=\mathrm{i} s_3$, these intervals of the imaginary axis are critical trajectories. Since elsewhere on the imaginary axis we have ${f(\eta,y,c)=h_\eta(\eta,y,c)^2<0}$, $\operatorname{Re}(h(\eta,y,c))$ is strictly monotone as $\eta$ varies in these intervals of $\mathrm{i}\mathbb{R}$. Therefore, in this configuration there are no points of either $K$ or $K'\subset K$ on the imaginary axis outside the two critical trajectories. Since $K'$ is symmetric in reflection through the imaginary axis, and since exactly one critical trajectory goes to $\eta=\infty$ in each half-plane, there are only three possibilities:
\begin{itemize}\itemsep=0pt
\item The two remaining trajectories emanating from $\eta=\mathrm{i} s_1$ tend to infinity in opposite half-planes, and there is a symmetric pair of arcs in each half-plane joining the points $\eta=\mathrm{i} s_2$ and $\eta=\mathrm{i} s_3$. However, since $\operatorname{Re}(h(\eta,y,c))$ is harmonic between the imaginary axis and each of these arcs and vanishes on each critical trajectory, this would imply by the maximum principle that $\operatorname{Re}(h(\eta,y,c))\equiv 0$ in each of these domains. This is a contradiction since $h_\eta(\eta,y,c)$ does not vanish identically.
\item The two remaining trajectories emanating from $\eta=\mathrm{i} s_2$ tend to infinity in opposite half-planes, and there is a symmetric pair of arcs in each half-plane joining the points $\eta=\mathrm{i} s_1$ and $\eta=\mathrm{i} s_3$. However, this would imply a crossing of two different trajectories at a point in each half-plane where $f(\eta,y,c)$ is finite and nonzero, which cannot occur.
\item Therefore, the remaining possibility must hold, namely that the two remaining trajectories emanating from $\eta=\mathrm{i} s_3$ tend to infinity in opposite half-planes, and there is a symmetric pair of arcs in each half-plane joining the points $\eta=\mathrm{i} s_1$ and $\eta=\mathrm{i} s_2$.
\end{itemize}
This shows that the claimed structure holds in case (i) when $s_1<0<s_2<s_3$.

If instead $s_1<0$ and $s_3=s_2^*$, since $f(\eta,y,c)=h_\eta(\eta,y,c)^2>0$ holds for $\eta\in \mathrm{i} \mathbb{R}$ between $\eta=\mathrm{i} s_1$ and $\eta=0$, this interval of the imaginary axis is a critical trajectory, while in the intervals between $\eta=-\mathrm{i}\infty$ and $\eta=\mathrm{i} s_1$ and between $\eta=0$ and $\eta=+\mathrm{i}\infty$ we have $f(\eta,y,c)=h_\eta(\eta,y,c)^2<0$ so $\operatorname{Re}(h(\eta,y,c))$ is strictly monotone. This implies that there are no points of either $K$ or $K'\subset K$ on the imaginary axis below $\eta=\mathrm{i} s_1$, but because $\operatorname{Re}(h(\eta,y,c))$ necessarily changes sign on the positive imaginary axis due to the singularity at the origin and the linear growth at infinity there is exactly one point of $K$ there, which may belong to $K'$. In fact, this point does indeed belong to $K'$, because otherwise at least two of the critical trajectories emanating from each of the points $\eta=\mathrm{i} s_2$ and $\eta=\mathrm{i} s_3=-(\mathrm{i} s_2)^*$ must tend to infinity in the half-plane containing the point because only one of them can terminate at $\eta=\mathrm{i} s_1$; this contradicts the fact that exactly one critical trajectory tends to infinity in each half-plane. So the distinguished point in the imaginary interval between $\eta=0$ and $\eta=+\mathrm{i}\infty$ belongs to $K'$ and lies on a critical trajectory crossing the imaginary axis horizontally and connecting $\eta=\mathrm{i} s_2$ and $\eta=\mathrm{i} s_3=-(\mathrm{i} s_2)^*$. The remaining two trajectories emanating from each of these points necessarily tend to $\eta=\mathrm{i} s_1$ and $\eta=\infty$ without crossing. This shows that the claimed structure of $K'$ holds in case (ii) when $s_1<0$ and $s_3=s_2^*$.

We next show that whenever $0<y<y_\mathrm{c}\approx 0.29177$, where $y_\mathrm{c}$ is the critical value defined in~\cite[Section 4.6]{BuckinghamM22}, there exists a unique $c=c_1(y)\in\mathbb{R}$ for which the conditions~\eqref{eq:Boutroux} (really just $I_{1,2}=0$ as $I_{2,3}=0$ is automatic) hold with root configuration $s_1<0<s_2<s_3$ and hence $K'$ has the claimed structure in case (i). To do this, we first suppose that $y>0$ and choose $c\in\mathbb{R}$ differently, so that the cubic $\mu\mapsto P(\mu,y,c)$ defined in~\eqref{eq:h-eta} has a simple root $\mu=s$ and a double root $\mu=d$.
Then by setting $P(\mu,y,c)=-(\mu-d)^2(\mu-s)$ one sees that $d=(1-s)/2$ and that $s$ satisfies the cubic equation $s(s-1)^2=-y^2$, while $c=-d^2-2ds=\frac{1}{4}(3s^2-2s-1)$. The condition $y>0$ implies that the equation $s(s-1)^2=-y^2$ has one real and two complex-conjugate solutions for $s$. But if $s=u+\mathrm{i} v$ with $v\neq 0$, then $\operatorname{Im}(c)=\frac{1}{2}(3u-1)v$ which vanishes for $v\neq 0$ only if $u=\frac{1}{3}$. Then $\operatorname{Im}\bigl(y^2\bigr)=-\operatorname{Im}\bigl(s(s-1)^2\bigr)=v^3\neq 0$ which contradicts $y>0$. Therefore, the conditions $y>0$ and $c\in\mathbb{R}$ require that we select the real root $s=s(y)<0$ of $s(s-1)^2=-y^2$ and then $d=d(y)>\frac{1}{2}$ is also real and $c=c_0(y)$ is a corresponding well-defined real number. In this double-root configuration, the function $\eta\mapsto h_\eta(\eta,y,c_0(y))$ is analytic except on the imaginary segment between $\eta=\mathrm{i} s(y)$ and $\eta=0$. Choosing an integration constant so that $\operatorname{Re}(h(\eta,y,c_0(y)))=0$ for $\eta=\mathrm{i} s(y)$, the function $\operatorname{Re}(h(\eta,y,c_0(y)))$ is well defined by contour integration and it is harmonic except on the branch cut for $h_\eta(\eta,y,c_0(y))$. According to~\cite[Section 4.6]{BuckinghamM22}, if $y>0$ and $y\neq y_\mathrm{c}$, then $y-y_\mathrm{c}$ and $\operatorname{Re}(h(\eta,y,c_0(y)))$ always have the same sign for $\eta=\mathrm{i} d(y)$. Now $c=c_0(y)$ is necessarily a real root of the cubic discriminant of $\mu\mapsto P(\mu,y,c)$, which is proportional by a positive numerical factor to $64 c^3+16c^2+72y^2c + 16y^2-27y^4$. The discriminant of this latter polynomial with respect to $c$ is proportional by a positive numerical factor to $-y^2\bigl(4+27y^2\bigr)^3$ which is strictly negative for $y>0$, hence $c=c_0(y)$ is a simple root of the cubic discriminant and it is the only real root thereof. Clearly the cubic discriminant of $\mu\mapsto P(\mu,y,c)$ has the same sign as $c$ when $y>0$ is fixed and $c\in\mathbb{R}$ is large. Hence $\mu\mapsto P(\mu,y,c)$ has three real distinct roots whenever $c>c_0(y)$ and has only one real root (simple) whenever~${c<c_0(y)}$.

Now returning to the general case of $y>0$ and $c\in\mathbb{R}$ arbitrary (so that $P(\mu,y,c)$ need not have a double root), we wish to solve the equation $I_{1,2}=0$ for $c$ using the above information. A~simple calculation shows that $\frac{\partial}{\partial c}h_\eta(\eta,y,c)^2 = -\eta^{-2}$, which implies that when $c\in\mathbb{R}$ and~${y>0}$, \looseness=1
\begin{equation*}
\frac{\partial I_{1,2}}{\partial c}=\operatorname{Re}\bigg(\int_{\mathrm{i} s_1}^{\mathrm{i} s_2}\frac{\partial^2 h}{\partial\eta\partial c}(\eta,y,c)\,\mathrm{d} y\bigg)=-\frac{1}{2}\operatorname{Re}\bigg(\int_{\mathrm{i} s_1}^{\mathrm{i} s_2}\frac{\mathrm{d}\eta}{\eta^2 h_\eta(\eta,y,c)}\bigg),
\end{equation*}
because $h_\eta(\eta,y,c)$ vanishes at $\eta=\mathrm{i} s_1$ and $\eta=\mathrm{i} s_2$. Some contour deformations show that regardless of whether $c<c_0(y)$ or $c>c_0(y)$, this derivative can be written in a universal form:
\begin{equation*}
\frac{\partial I_{1,2}}{\partial c}=\frac{1}{4}\int_\mathbb{R}\frac{\chi_{P>0}(\mu)\,\mathrm{d}\mu}{\sqrt{P(\mu,y,c)}},
\end{equation*}
where $\chi_{P>0}(\mu)$ is the characteristic function of the union of intervals of $\mathbb{R}$ on which $\mu\mapsto P(\mu,y,c)$ is positive, and the square root is positive. Therefore, in both cases $c<c_0(y)$ and $c>c_0(y)$ the partial derivative of $I_{1,2}$ with respect to $c\in\mathbb{R}$ for $y>0$ is positive (the difference is that the support of $\chi_{P>0}(\mu)$ consists of two intervals $(-\infty,s_1)\cup (0,+\infty)$ for $c<c_0(y)$ and of three intervals $(-\infty,s_1)\cup (0,s_2)\cup (s_3,+\infty)$ for $c>c_0(y)$). Now, the real-valued function $c\mapsto I(c):=I_{1,2}$ is certainly continuous as a function of $c\in\mathbb{R}$ and is continuously differentiable for $c\neq c_0(y)$ with positive derivative. We also know that $I(c_0(y))=\operatorname{Re}(h(\mathrm{i} d(y),y,c_0(y)))<0$ for $0<y<y_\mathrm{c}$. However, it is also true that eventually $I(c)>0$ as $c\to+\infty$. To see this, first note that the roots of $P(\mu,y,c)$ for large $c>0$ with fixed $y>0$ are $s_1=-c^{1/2}+O(1)$, $s_2=\frac{1}{4}y^2c^{-1}+O\bigl(c^{-2}\bigr)$, and $s_3=c^{1/2}+O(1)$. Then one rescales the integrand of $I_{1,2}$ by $\eta=s_2w$ and notices that by contour deformations, the leading term of $I_{1,2}$ as $c\to+\infty$ with $y>0$ fixed is computed as a residue. Indeed, letting $C_\mathrm{L}$ (resp.\ $C_\mathrm{R}$) denote a contour in the left (resp.\ right) half-plane beginning at a point on the imaginary axis between $\eta=\mathrm{i} s_1$ and $\eta=0$ and terminating at a point on the imaginary axis between $\eta=\mathrm{i} s_2$ and $\eta=\mathrm{i} s_3$, letting $L$ denote a~positively-oriented loop surrounding the imaginary interval between $w=0$ and $w=\mathrm{i}$, and using principal branches for all power functions,
\begin{align*}
I_{1,2}&=\operatorname{Re}\left(\int_{\mathrm{i} s_1}^{\mathrm{i} s_2}h_\eta(\eta,y,c)\,\mathrm{d}\eta\right) \\ &=
\frac{1}{2}\operatorname{Re}\left(\int_{C_\mathrm{L}\cup C_\mathrm{R}}\left[-\mathrm{i}(-\mathrm{i}\eta-s_3)^{1/2}(-\mathrm{i}\eta-s_2)^{1/2}(-\mathrm{i}\eta)^{-3/2}(-\mathrm{i}\eta-s_1)^{1/2}\right]\,\mathrm{d}\eta\right)\\
&=\frac{1}{2}\operatorname{Re}\left(\oint_{(-C_\mathrm{L})\cup C_\mathrm{R}}\left[-\mathrm{i} (-\mathrm{i}\eta-s_1)^{1/2}\left(\frac{\eta-\mathrm{i} s_2}{\eta}\right)^{1/2}(\mathrm{i}\eta+s_3)^{1/2}\right]\frac{\mathrm{d}\eta}{\eta}\right)\\
&=\frac{1}{2}\sqrt{c}\operatorname{Re}\left(\oint_{L}\left[-\mathrm{i}\left(-\mathrm{i}\frac{s_2}{\sqrt{c}}w-\frac{s_1}{\sqrt{c}}\right)^{1/2}\left(\frac{w-\mathrm{i}}{w}\right)^{1/2}\left(\mathrm{i}\frac{s_2}{\sqrt{c}}+\frac{s_3}{\sqrt{c}}\right)^{1/2}\right]\,\frac{\mathrm{d} w}{w}\right)\\
&=\frac{1}{2}\sqrt{c}\left[\operatorname{Re}\left(\oint_L\left[-\mathrm{i}\left(\frac{w-\mathrm{i}}{w}\right)^{1/2}\right]\frac{\mathrm{d} w}{w}\right)+o(1)\right] \\ %,\quad c\to+\infty\\
&=\frac{1}{2}\sqrt{c}\left[\operatorname{Re}\left(2\pi\mathrm{i}\mathop{\mathrm{Res}}_{w=\infty}
\left[-\mathrm{i}\left(\frac{w-\mathrm{i}}{w}\right)^{1/2}\frac{1}{w}\right]\right)+o(1)\right] \\ %,\quad c\to+\infty\\
&=\pi\sqrt{c}+o(\sqrt{c}),\qquad c\to+\infty.
\end{align*}
Therefore if $0<y<y_\mathrm{c}$, $I(c)<0$ for $c=c_0(y)$ while $I(c)>0$ in the limit $c\to+\infty$, and $c\mapsto I(c)$ is strictly increasing for $c>c_0(y)$. It follows from the intermediate value theorem that there is a unique solution $c=c_1(y)>c_0(y)$ of $I(c)=I_{1,2}=0$, and the inequality $c_1(y)>c_0(y)$ implies that the roots of $\mu\mapsto P(\mu,y,c)$ are all real and hence $K'$ has the claimed structure in case (i) whenever $0<y<y_\mathrm{c}$.

Now we continue this solution into the complex $y$-plane. Since for $0<y<y_\mathrm{c}$ the roots of $\mu\mapsto P(\mu,y,c_1(y))$ are distinct, the solution $c=c_1(y)$ of the system~\eqref{eq:Boutroux} can be continued uniquely by the implicit function theorem to some maximal domain of the complex $y$-plane containing the real interval $(0,y_\mathrm{c})$ (however note that $c_1(y)$ is not an analytic function of $y$ in this domain). The boundary of this domain consists of all points $y$ for which $c_1(y)=c_0(y)$, i.e., under continuation of the solution of the Boutroux conditions~\eqref{eq:Boutroux} one arrives at a degenerate spectral curve with $\mu\mapsto P(\mu,y,c)$ having a double root. This boundary was obtained in~\cite[Section 4.7]{BuckinghamM22}, which shows that the maximal domain under consideration is exactly that mapped by $Y=y^{1/3}$ onto the interior of the right ``wing'' of the ``bow-tie'' region $\mathcal{B}$ in the $Y$-plane (see Figure~\ref{fig:DensityPlots}, left-hand panel). The topological structure of $K'$ remains the same under continuation, and hence the claimed structure of $K'$ holds throughout this domain in case (i) as that is the case for $0<y<y_\mathrm{c}$. The boundary of the domain consists of the imaginary segment between $y=\pm 2\mathrm{i}/
\sqrt{27}$ and a certain Schwarz-symmetric curve in the right half-plane connecting those two imaginary endpoints via the positive real value $y_\mathrm{c}\approx 0.29177$; see the right-hand panel of Figure~\ref{fig:DensityPlots}.\looseness=-1

Given the proven structure of $K'$ in case~(i), it is clear that $K'$ divides the complex $\eta$-plane into three disjoint regions, consistent with the \emph{basic structure theorem}~\cite[p.~37]{Jenkins58}. As each of these regions has exactly one pole of $f(\eta,y,c)\,\mathrm{d}\eta^2$ on its boundary (including the point at infinity in a suitable local coordinate), they are all \emph{end domains}, which means that they are conformally mapped by the primitive $\eta\mapsto \int^\eta h_\eta(\eta',y,c)\,\mathrm{d}\eta'$ onto a half-plane with a~vertical boundary. Since $\operatorname{Re}(h(\eta,y,c))=0$ on the boundary of each end domain because it is a subset of $K'\subset K$, the vertical boundary of the image is exactly the imaginary axis, and hence there can be no other points with $\operatorname{Re}(h(\eta,y,c))=0$ in the interior of each end domain. This proves that $K=K'$, which establishes the key properties of the level curve $K$ asserted in the statement of the lemma.

Introducing a contour arc $\Sigma_{0,2}$ joining $\eta=0$ with $\eta=\mathrm{i} s_2$ and an unbounded arc $\Sigma_{1,\infty}$ with finite endpoint $\eta=\mathrm{i} s_1$ such that $\Sigma_h:=\Sigma_{2,3}\cup\Sigma_{0,2}\cup\Sigma_{0,1}\cup\Sigma_{1,\infty}$ is a simple contour, a function $\eta\mapsto h(\eta,y,c_1(y))$ is well-defined up to an integration constant possibly depending on $y$ by contour integration of $h_\eta(\eta,y,c_1(y))$ in the simply connected domain~${\mathbb{C}\setminus\Sigma_h}$. Since $h_\eta(\eta,y,c)$ is integrable at all three of the roots $\eta=\mathrm{i} s_j$, we can and will choose the integration constant so that $\operatorname{Re}(h(\mathrm{i} s_3,y,c_1(y)))=0$. It follows from~\eqref{eq:Boutroux} and the reality of the residue of $h_\eta(\eta,y,c_1(y))$ at $\eta=\infty$ that $\operatorname{Re}(h(\eta,y,c_1(y)))$ is a function harmonic on the larger domain $\mathbb{C}\setminus(\Sigma_{2,3}\cup\Sigma_{0,1})$. Recalling that $\Sigma_{2,3}$ and $\Sigma_{0,1}$ have been chosen to agree~with arcs of $K'=K$, the function $\eta\mapsto\operatorname{Re}(h(\eta,y,c_1(y)))$ is continuous except at ${\eta=0}$.
\end{proof}
